\chapter{Dual abstract-concrete methodology}
\label{app:dual-methodology}


\section{Methodology}
\label{sec:philosophy}

\begin{convention}[Dual proof methodology]
\label{conv:dual-proof}
Several results in this monograph admit two independent proofs: one via abstract $\infty$-categorical machinery (Verdier duality, operadic Koszul theory, Francis--Gaitsgory) and one via explicit chain-level constructions (bar complex differentials, spectral sequences, contracting homotopies). We present both where they illuminate different aspects of the same phenomenon: the abstract proof establishes existence and functoriality, while the concrete proof provides computable formulas.
\end{convention}

\begin{remark}[Dual proofs]
\label{rem:dual-proofs}
Agreement between abstract and concrete approaches provides internal validation: errors in one approach are typically detected by the other. Abstract proofs generalize more readily to derived settings; concrete proofs reveal structures visible only with explicit formulas.
\end{remark}


\section{Key instances: bar-cobar and Riemann--Hilbert}
\label{sec:key-instances}

\subsection{Bar-cobar equivalence}

\begin{theorem}[Abstract bar-cobar equivalence \cite{FG12, HA}; \ClaimStatusProvedElsewhere]
\label{thm:abstract-bar-cobar}
In the $\infty$-category $\DMod^{\mathrm{fact}}(\Ran(X))$ of factorization
D-modules on the Ran space of a curve $X$, equipped with the chiral tensor
structure:
\[
\Bbar: \Eone\text{-}\Alg \rightleftarrows \Eone\text{-}\CoAlg : \Cobar
\]
is an adjoint equivalence of $\infty$-categories.
\end{theorem}

\begin{proof}[Proof (following Francis--Gaitsgory \cite{FG12})]
The chiral tensor structure on $\DMod(\Ran(X))$ is pro-nilpotent in the sense
of Francis--Gaitsgory: for any object $\cM$, the tensor powers $\cM^{\otimes n}$
supported on $\Ran_{\geq n}(X)$ have vanishing homology as $n \to \infty$
in the appropriate sense.

By the general theory of Koszul duality for operads in pro-nilpotent tensor
$\infty$-categories~\cite{FG12}, the bar-cobar adjunction is an equivalence.
Pro-nilpotence ensures the bar and cobar constructions
are completed and the unit/counit maps are equivalences.
\end{proof}

\begin{theorem}[Concrete bar-cobar equivalence \cite{LV12}; \ClaimStatusProvedElsewhere]
\label{thm:concrete-bar-cobar}
For an $\Eone$-chiral algebra $\cA$, the counit:
\[
\varepsilon: \Cobar^{\mathrm{geom}}(\Bbar^{\mathrm{geom}}(\cA)) \xrightarrow{\simeq} \cA
\]
is a quasi-isomorphism.
\end{theorem}

\begin{proof}
\emph{Step 1.} The double bar-cobar complex $\Cobar(\Bbar(\cA))$ has a
bigrading by cobar wordlength~$q$ (tensor degree in $T(s^{-1}\overline{B}(\cA))$)
and bar degree~$p$ (within each bar element). Filter by cobar wordlength~$q$.

\emph{Step 2.} The $E_0$-page at wordlength~$q$ is $(s^{-1}\overline{B}(\cA))^{\otimes q}$.
The $d_0$-differential is the bar differential~$d_{\bar{B}}$ (acting on each tensor factor,
preserving wordlength).

\emph{Step 3.} The \emph{augmented} (two-sided) bar complex
$B(\cA, \cA, \bC) = \cA \otimes T(s^{-1}\bar{\cA}) \otimes \bC$
admits a contracting homotopy $h[a_1|\cdots|a_n] = [1|a_1|\cdots|a_n]$
(the extra degeneracy) and is therefore acyclic. Note: this contracting
homotopy does \emph{not} restrict to the reduced bar complex~$\Bbar(\cA) = T^c(s^{-1}\bar{\cA})$
(since inserting~$1$ leaves~$\bar{\cA}$), so $\Bbar(\cA)$ itself is generally
\emph{not} acyclic; its homology computes~$\mathrm{Tor}^{\cA}_*(\bC, \bC)$.

\emph{Step 4.} We apply the \emph{comparison lemma} for twisted tensor
products (\cite{LV12}~\S2.1). Let $\tau: \Bbar(\cA) \to \cA$
be the universal twisting morphism (projection to the cogenerators) and
$\kappa: \Bbar(\cA) \to \Cobar(\Bbar(\cA))$ the canonical twisting
morphism (inclusion of cogenerators). By Step~3,
$\cA \otimes_\tau \Bbar(\cA) = B(\cA, \cA, \bC)$ is acyclic.
The twisted tensor product
$\Cobar(\Bbar(\cA)) \otimes_\kappa \Bbar(\cA) = B(\Cobar(\Bbar(\cA)),
\Cobar(\Bbar(\cA)), \bC)$ is likewise acyclic by the same extra-degeneracy
argument applied to the augmented algebra $\Cobar(\Bbar(\cA))$.
Both are acyclic twisted tensor products over the same coalgebra
$\Bbar(\cA)$, so the comparison lemma gives a quasi-isomorphism
$\Cobar(\Bbar(\cA)) \xrightarrow{\simeq} \cA$.
See Theorem~\ref{thm:bar-cobar-isomorphism-main} for the chiral generalization.
\end{proof}

\subsection{Riemann--Hilbert correspondence}

\begin{theorem}[Abstract Riemann--Hilbert \cite{KS90}; \ClaimStatusProvedElsewhere]
\label{thm:abstract-rh}
There is an equivalence of $\infty$-categories:
\[
\mathrm{RH}: \DMod(X)_{\mathrm{hol}} \xrightarrow{\simeq} \mathrm{Shv}_c(X^{\mathrm{an}}; k)
\]
between holonomic D-modules on a smooth variety $X$ and constructible sheaves
on the analytification $X^{\mathrm{an}}$. This equivalence:
\begin{enumerate}[label=(\roman*)]
\item Is functorial with respect to proper pushforward and $!$-pullback.
\item Is monoidal with respect to $\otimes^!$ and the derived tensor product.
\item Intertwines Verdier duality on both sides.
\end{enumerate}
\end{theorem}

\begin{proof}[Abstract proof]
The Riemann--Hilbert correspondence is characterized uniquely by the conditions
above: functoriality, monoidality, and compatibility with duality. Existence follows
from the general theory of regular holonomic D-modules: the de Rham functor
$\DR: \DMod_{\mathrm{hol}} \to \mathrm{Perv}(X^{\mathrm{an}})$ extends to an
equivalence, and perverse sheaves embed fully faithfully in constructible sheaves.
\end{proof}

\begin{construction}[Concrete Riemann--Hilbert]
\label{constr:concrete-rh}
For a holonomic D-module $\cM$ on $\FM_n(X)$ with regular singularities along
the boundary $D = \partial \FM_n(X)$:

\emph{Step 1.} On the open $\Conf_n(X) = \FM_n(X) \setminus D$, the D-module
$\cM$ is a vector bundle $\cV$ with flat connection $\nabla$.

\emph{Step 2.} The flat sections of $\nabla$ form a local system:
\[
\cL := \ker(\nabla: \cV \to \cV \otimes \Omega^1) \to \Conf_n(X).
\]

\emph{Step 3.} Near a boundary stratum $D_T$ (labeled by
a tree~$T$), choose local coordinates
$(z_1, \ldots, z_n,\allowbreak t_1, \ldots, t_k)$
with $D_T = \{t_1 \cdots t_k = 0\}$. The connection has the
form:
\[
\nabla = d + \sum_{i=1}^k A_i \frac{dt_i}{t_i} + \text{(regular terms)}
\]
with $A_i \in \End(\cV)$ the \emph{residue matrices}.

\emph{Step 4.} Regular singularity means the monodromy around $\{t_i = 0\}$ is:
\[
M_i = \exp(-2\pi i A_i)
\]
which is quasi-unipotent (the semisimple part of $M_i$ has finite order) for D-modules with regular singularities and rational exponents.

\emph{Step 5.} The Riemann--Hilbert sheaf $\mathrm{RH}(\cM)$ is $\cL$ extended
to $\FM_n(X)$ as a constructible sheaf by taking nearby cycles along $D$.
\end{construction}


\subsection{\texorpdfstring{$\infty$-operads}{infinity-operads}}\label{sec:infty-operads}

\begin{definition}[\texorpdfstring{$\infty$}{infinity}-operad]
\label{def:infty-operad}
An \emph{$\infty$-operad} is a functor $\cO^\otimes \to \mathrm{Fin}_*$
satisfying the Segal conditions (Lurie, HA Definition~2.1.1.10), where $\mathrm{Fin}_*$ is the category of
pointed finite sets. A \emph{symmetric monoidal $\infty$-category} is the special case where $\cO^\otimes \to \mathrm{Fin}_*$ is also a coCartesian fibration; this is strictly more restrictive than the general notion of $\infty$-operad.
\end{definition}

\begin{theorem}[Geometric models for \texorpdfstring{$\infty$}{infinity}-operads; \ClaimStatusProvedHere]
\label{thm:geometric-infty-operads}
The geometric bar and cobar constructions provide explicit models for the
$\infty$-operadic bar-cobar theory:
\begin{enumerate}[label=(\roman*)]
\item The bar complex $\Bbar^{\mathrm{geom}}(\cA)$ is a model for the
$\infty$-categorical bar construction $\Bbar_\infty(\cA)$.

\item The cobar complex $\Cobar^{\mathrm{geom}}(\cC)$ is a model for the
$\infty$-categorical cobar construction $\Cobar_\infty(\cC)$.

\item Quasi-isomorphisms of geometric models correspond to equivalences of
$\infty$-categorical objects.
\end{enumerate}
\end{theorem}

\begin{proof}
The geometric constructions are ``strictly'' models for the
derived constructions:

\emph{For bar.} The geometric bar complex $\Bbar^{\mathrm{geom}}(\cA)$ is
computed as sections of a factorization algebra on configuration spaces.
By factorization homology (Ayala--Francis), this computes the $\infty$-categorical
tensor product:
\[
\Bbar^{\mathrm{geom}}_n(\cA) = \int_{\FM_n(X)} \cA^{\boxtimes n} = \cA^{\otimes^!_\infty n}
\]
where $\otimes^!_\infty$ is the derived $!$-tensor product of factorization
D-modules.

\emph{For cobar.} Similarly, the distributional cobar complex computes
the $\infty$-categorical cofree coalgebra on the suspended cogenerators.
\end{proof}

\begin{corollary}[Bar-cobar inversion establishes Koszulness \cite{LV12}; \ClaimStatusConditional]
\label{cor:derived-koszul}
\emph{Type signature:} (Open quadrant; bar--cobar; Beilinson level
$1\!\leftrightarrow\!2$; HP = augmented + complete + Koszul locus).
For any augmented $\Eone$-chiral algebra $\cA$ on the chirally Koszul
locus, the counit of the bar-cobar adjunction
\[
\varepsilon\colon \Cobar_\infty(\Bbar_\infty(\cA)) \xrightarrow{\simeq} \cA
\]
is an equivalence in the $\infty$-category of $\Eone$-chiral algebras. This
is the counit face of Theorem~A (bar-cobar inversion), not a definition of
Koszulness: the
equivalence $\Omega B(\cA) \simeq \cA$ holds unconditionally on the chirally
Koszul locus, and characterizes it (see the Koszul equivalences meta-theorem,
Theorem~\ref{thm:koszul-equivalences-meta}).

\noindent\emph{Three functors, one input}: $B(\cA)$ is a
factorization coalgebra; $\Omega(B(\cA)) \simeq \cA$ recovers the original
algebra (inversion); $\mathbb{D}_{\mathrm{Ran}}(B(\cA)) \simeq B(\cA^!)$
produces the Koszul dual bar coalgebra (Verdier duality, Theorem~A).
These are distinct operations.
\begin{proof}
Theorem~\ref{thm:concrete-bar-cobar} gives the comparison quasi-isomorphism at
the chain level; Theorem~\ref{thm:abstract-bar-cobar} upgrades it to an
$\infty$-categorical equivalence.
\end{proof}
\end{corollary}

\subsection{Gui--Li--Zeng quadratic duality: candidate dual with MC
comparison map, not full Koszulness}
\label{subsec:glz-scope-vs-full-koszulness}

The Gui--Li--Zeng quadratic duality~\cite{GLZ22} produces a
\emph{candidate} Koszul dual with an associated Maurer--Cartan
comparison map. Full Koszulness — the bar--cobar inversion
$\Omega B(\cA) \simeq \cA$ on the locus — is a strictly stronger
statement that requires three independent ingredients. This
distinction is load-bearing and is the MA-12 scope discipline applied
to quadratic chiral duality.

\begin{theorem}[Gui--Li--Zeng quadratic duality: scope; \ClaimStatusProvedElsewhere]
\label{thm:glz-quadratic-duality-scope}
\emph{Type signature:} (Open quadrant; quadratic chiral data;
Beilinson level $1\!\rightarrow\!2$; HP = $E_\infty$-chiral +
quadratic presentation $(N,P)$ + dualizable quadratic datum
$(s^{-1}N^\vee\omega^{-1}, P^\perp)$).
Let $\cA = \cA(N,P)$ be a quadratic chiral algebra with dualizable
quadratic datum in the sense of~\cite{GLZ22}. Then:
\begin{enumerate}[label=\textup{(\arabic*)}]
\item \textbf{Candidate quadratic dual.} There is a quadratic chiral
algebra
\[
\cA^!_{\mathrm{GLZ}} := \cA\bigl(s^{-1}N^\vee\omega^{-1}, P^\perp\bigr),
\]
the GLZ quadratic dual.

\item \textbf{Maurer--Cartan comparison map.} There is a bijection
\[
\operatorname{Hom}_{\mathrm{ChirAlg}}(\cA, \cB)
\;\longleftrightarrow\;
\MC\bigl(\cA^!_{\mathrm{GLZ}} \otimes \cB\bigr),
\]
between chiral algebra morphisms and Maurer--Cartan elements in the
tensor product convolution algebra.
\end{enumerate}
The output is a candidate dual together with a comparison map, not a
witnessed bar--cobar inversion.
\end{theorem}

\begin{proof}
This is~\cite[Proposition~3.2 and Theorem~5.3]{GLZ22}, restated
verbatim with type signature attached.
\end{proof}

\begin{theorem}[Full Koszulness as a separate theorem; \ClaimStatusConditional]
\label{thm:full-koszulness-separate}
\emph{Type signature:} (Open quadrant; bar--cobar inversion;
Beilinson level $1\!\leftrightarrow\!2$; HP = (a)+(b)+(c) below).
Let $\cA = \cA(N,P)$ be a quadratic chiral algebra. Then $\cA$ is
chirally Koszul — meaning the bar--cobar inversion
$\Cobar^{\mathrm{ch}}\Bbar^{\mathrm{ch}}(\cA) \xrightarrow{\sim} \cA$
holds in the appropriate ambient — \emph{if and only if} all three
of the following independent ingredients hold simultaneously:
\begin{enumerate}[label=\textup{(\alph*)}]
\item \textbf{Quadratic-dual candidate exists.} The dualizable
quadratic datum $(s^{-1}N^\vee\omega^{-1}, P^\perp)$ is well-formed
and produces $\cA^!_{\mathrm{GLZ}}$
(Theorem~\ref{thm:glz-quadratic-duality-scope}).

\item \textbf{Maurer--Cartan comparison map is a quasi-isomorphism.}
The comparison map of
Theorem~\ref{thm:glz-quadratic-duality-scope}(2) lifts to a
quasi-isomorphism of convolution dg Lie algebras
$\operatorname{Conv}(\cA^{\mathrm i}, \cA^!_{\mathrm{GLZ}}) \to
\operatorname{Conv}(\cA^{\mathrm i}, \cA)$ identifying the
canonical twisting morphism $\tau$ with the GLZ Maurer--Cartan
element $\alpha$.

\item \textbf{Bar--cobar adjunction $K^2 \simeq \id$ on the comparison
locus.} The counit $\Cobar^{\mathrm{ch}}\Bbar^{\mathrm{ch}}(\cA)
\xrightarrow{\sim} \cA$ is a chain-level quasi-isomorphism on the
locus where (a) and (b) hold; equivalently, the PBW spectral sequence
degenerates at $E_2$ on this locus.
\end{enumerate}
Each ingredient is independent: (a) without (b) gives a candidate
without comparison; (a)+(b) without (c) gives a Maurer--Cartan
comparison without bar--cobar inversion (the higher
$A_\infty$-operations $m_k$, $k\ge 3$, may obstruct
strictification); (a)+(c) without (b) is impossible (the comparison
map is forced by the quadratic datum once dualisability is granted).
\end{theorem}

\begin{proof}
\textbf{Necessity.} If bar--cobar inversion holds, the bar coalgebra
is concentrated on quadratic generators and quadratic relations
(degrees~$1$ and~$2$ in bar weight), forcing the dualizable quadratic
datum and the comparison map. The $E_2$-degeneration of the PBW
spectral sequence is the chain-level witness.

\textbf{Sufficiency.} Granted (a)+(b)+(c), the comparison map of (b)
identifies $\Cobar^{\mathrm{ch}}(\Bbar^{\mathrm{ch}}(\cA))$ with the
quadratic algebra reconstructed from the GLZ dual on the locus
of~(c); the $E_2$-degeneration of (c) ensures no higher $A_\infty$
shadows obstruct the inversion.
\end{proof}

\begin{remark}[Why all three ingredients are required: MA-12]
\label{rem:glz-three-ingredients-ma12}
\index{MA-12 (finite-spin / quadratic / classical = full theorem)}
Asserting GLZ quadratic duality ``establishes Koszulness'' without
ingredients (b) and (c) is a textbook MA-12 violation: a
quadratic-only construction is promoted to the full Koszul-duality
theorem without naming the endpoint hypothesis package. The endpoint
hypotheses are (b) the Maurer--Cartan comparison map being a
quasi-isomorphism and (c) the bar--cobar adjunction $K^2 \simeq \id$
on the comparison locus, which specialise the
Vallette~\cite{Val16} / Hinich~\cite{Hinich2003} /
Positselski~\cite{Positselski11,Positselski2018} bar--cobar Quillen
package to the chiral ambient.

The MA-12 violation is detected by the slogan \emph{``GLZ produces
the Koszul dual''} — the truth is GLZ produces a \emph{candidate}
dual with a comparison map; the inversion theorem is separate. The
companion-theorem package is named in the type signature of
Theorem~\ref{thm:full-koszulness-separate}.
\end{remark}

\begin{remark}[Companion: chain-level versus $\infty$-categorical lanes]
\label{rem:glz-two-lanes}
The three-ingredient package of
Theorem~\ref{thm:full-koszulness-separate} has two lanes,
parallelling the dual-proof methodology of
Convention~\ref{conv:dual-proof}.
\emph{Chain-level (Priddy lane).} Ingredient (c) is witnessed by an
explicit Priddy--Koszul chain homotopy contracting the augmentation
ideal, with $E_2$-degeneration of the PBW spectral sequence
weight-by-weight (\cite[\S7.4]{LV12}; chiral specialisation).
\emph{$\infty$-categorical (Positselski lane).} Ingredient (c) is
witnessed by the coderived/contraderived bar--cobar Quillen
equivalence on the pro-conilpotent completion
$\widehat{\cA} = \varprojlim_N \cA/\cA_{\geq N}$
(\cite{Positselski2011}; chiral specialisation in
Theorem~\ref{thm:abstract-bar-cobar}).
The two lanes coincide on the Koszul locus and diverge off it; bare
GLZ quadratic duality alone does not select either lane.
\end{remark}
